\documentclass[a4paper]{article}

\usepackage[english]{babel}
%\usepackage[utf8x]{inputenc}
\usepackage{amsmath}
\usepackage{graphicx}
\usepackage[colorinlistoftodos]{todonotes}
\usepackage[all]{xy}
\usepackage{dsfont}
\usepackage{amssymb}
\usepackage{bbold}
\usepackage{amsthm}

\usepackage{tikz}
\usepackage{verbatim}
\usetikzlibrary{matrix}
\usetikzlibrary{cd}

\usepackage[autostyle=true]{csquotes}
\usepackage{pdfpages}
\usepackage{hyperref}
\usepackage{fancyhdr}



\newtheorem{theorem}{Theorem}[section]
\newtheorem{proposition}{Proposition}[theorem]
\newtheorem{lemma}{Lemma}[theorem]
\newtheorem{corollary}{Corollary}[theorem]
\newtheorem*{definition}{Definition}


\newcommand{\norm}[1]{\left\lVert#1\right\rVert}

\newcommand{\Rn}{\mathbb{R}^n}


\newcommand{\too}[1]{\overset{#1}{\longrightarrow}}
\newcommand{\hookar}[1]{\overset{#1}{\hookrightarrow}}
\newcommand{\colim}[1]{ \underset{#1}{\texttt{colim}} \text{ } }
\newcommand{\fset}[1]{ \{ #1_0 , \dots , #1_n \} }
\newcommand{\ddt}{\frac{d}{dt}}


\newcommand{\N}{\mathbb{N}}
\newcommand{\Z}{\mathbb{Z}}







\fancyhead[L]{\includegraphics[scale=0.1]{UFT.png}}

\begin{document}


\thispagestyle{fancy}
\vspace*{1cm}

\begin{minipage}[r]{12cm}
    \begin{flushright}
\begin{tikzpicture}
\draw[green, very thick] (0,0) rectangle (3,2);
\end{tikzpicture}
    \end{flushright}
\end{minipage}

\begin{center}
\Large{\textbf{Control 1\\ Álgebra Lineal\\ }}
\end{center}
\vspace{1cm}
% \begin{minipage}[r]{15cm}
% \begin{flushleft}

% Fecha: 15 de abril 2024 \\%
% Duración: 120 minutos  \\
% Profesores: Camila Muñoz.\\\hspace{1.7cm} Adrián Silva.\\
% \end{flushleft}
% \end{minipage}
% \vspace{2cm}
% \begin{center}
% \Large{\textbf{Facultad de Ingeniería}}
% \end{center}
% \vspace{1cm}
% \textbf{En caso de que un alumno sea sorprendido en actitudes deshonestas, ya sea intentando o
% dejándose copiar, será calificado con nota mínima 1.0 y perderá toda posibilidad de eliminar
% dicha nota. Además, esta conducta quedará archivada como antecedente en su carpeta y
% podrá ser causal de eliminación.
% Así mismo, está prohibido salir de la sala, usar o mantener sobre la mesa teléfonos celulares
% o cualquier tipo de aparato electrónico, a no ser que el profesor lo requiera.}\\



% \vspace{5cm}
% \makebox[12cm]{\textbf{Nombre Alumno}:\ \hrulefill} 
% \newpage
\thispagestyle{empty}
\section*{Preguntas}
\begin{enumerate}


    \item {\bf (XX puntos)} Usando el teorema del rango, determine los valores de $\alpha$ de modo que el sistema
    \begin{align*}
 x + 3y &= 1 \\
 x + 2y +z &= 0\\
2 x + 5y +(\alpha -1) z &= 1-\alpha 
\end{align*}

    tenga solución única, infinitas soluciones o ninguna solución. Para $\alpha=0$, determine el conjunto soluci\'on del sistema.

        \item {\bf (XX puntos)} Usando el teorema del rango, determine los valores de $\alpha$ de modo que el sistema
    \begin{align*}
 x + y &= -\alpha \\
\alpha x + 2y +z &= 1\\
 x + y + z &= 1-\alpha 
\end{align*}

    tenga solución única, infinitas soluciones o ninguna solución. Para $\alpha=1$, determine el conjunto soluci\'on del sistema.
        
        









\end{enumerate}


\end{document}